\article{$\GL(2)$ 的局部 Langlands 猜想报告}{Report on the local Langlands conjecture for GL(2)}{J. Tunnell}{2}{135--138}

记 $\Phi(\GL(2)/K)$ 为非 Archimedes 局部域（nonarchimedean local field）$K$ 的 Weil--Deligne 群（Weil--Deligne group）$W'_K$ 的二维 $F$-半单表示（$F$-semisimple representation）的同构类集合。本报告讨论 Langlands 的下述猜想。这一猜想将 $\Phi(\GL(2)/K)$ 与 $\GL(2,K)$ 的不可约容许表示（irreducible admissible representation）的同构类集合 $\Pi(\GL(2,K))$ 联系起来。

\begin{originaltheorem}{猜想（Conjecture）}
对 $\Phi(\GL(2)/K)$ 中每个表示 $\sigma$，存在 $\Pi(\GL(2,K))$ 中的表示 $\pi=\pi(\sigma)$，使 $\sigma$ 的行列式与 $\pi$ 的中心拟特征（central quasi-character）相等，并且对 $K^*$ 的所有拟特征（quasi-character）$\chi$ 都有
\[
 L(\pi\otimes\chi)=L(\sigma\otimes\chi),
 \qquad \varepsilon(\pi\otimes\chi)=\varepsilon(\sigma\otimes\chi).
\]
映射 $\sigma\mapsto\pi(\sigma)$ 是 $\Phi(\GL(2)/K)$ 到 $\Pi(\GL(2,K))$ 的双射。
\end{originaltheorem}

猜想中的存在性断言由 [4] 提出。该文还证明，给定 $\sigma$ 后，满足上述条件的 $\pi$ 至多只有一个。单射性断言等价于说，$W'_K$ 的二维 $F$-半单表示由其扭曲的 $L$ 因子和 $\varepsilon$ 因子（twisted $L$- and $\varepsilon$-factors）以及行列式确定。可约二维表示仅凭各个扭曲 $L$ 因子即可确定，这一点可以直接证明。对于不可约表示，已知的单射性证明则借助容许表示理论。

如 [1, 3.2.3] 所述，由 Jacquet 与 Langlands 的工作可知，当 $\sigma$ 可约时 $\pi(\sigma)$ 存在，并由此建立了 $W'_K$ 的完全可约（completely reducible）（分别为可约但不可分解（reducible indecomposable））的二维 $F$-半单表示的同构类集合，与 $\GL(2,K)$ 的主系列表示（principal series representation）（分别为特殊表示（special representation））的同构类集合之间的双射。

令 $\Phi_{\mathrm{irr}}(\GL(2)/K)$ 为 Weil 群（Weil group）$W_K$ 的不可约二维表示的同构类集合，令 $\Pi_{\mathrm{cusp}}(\GL(2,K))$ 为 $\GL(2,K)$ 的超尖点表示（supercuspidal representation）的同构类集合。这两个子集都可用同一条件刻画：其中元素的所有扭曲 $L$ 因子均等于 $1$。为证明猜想，只需验证对每个 $\sigma\in\Phi_{\mathrm{irr}}(\GL(2)/K)$，$\pi(\sigma)$ 都存在，并证明由此建立了 $\Phi_{\mathrm{irr}}(\GL(2)/K)$ 与 $\Pi_{\mathrm{cusp}}(\GL(2,K))$ 之间的双射。这首先在剩余特征（residue characteristic）为奇数的局部域上得到证明（见第 2 节），并由此启发了猜想中一般情形的双射性断言。\footnote{原文的 AMS（MOS）主题分类（1970）：主要分类 12B30、12B15；次要分类 10D99。原文版权：1979，American Mathematical Society。}

除那些次数大于 $1$ 且不含三次单位根的 $\Q_2$ 的扩域外，这一猜想在其他所有非 Archimedes 局部域上都已得到证明。下文将综述这些结果，并列出相关证明的出处。

\origsection{1}{$\pi(\sigma)$ 的存在性（Existence of $\pi(\sigma)$）}
如果 $\sigma$ 是由 $W_K$ 的一个指数为二的子群的一维表示诱导而来的表示，则 Weil 的显式构造给出了满足要求的 $\GL(2,K)$ 的不可约容许表示 $\pi(\sigma)$ [4, 4.7]。当 $K$ 的剩余特征为奇数时，$\Phi_{\mathrm{irr}}(\GL(2)/K)$ 中的所有表示都由真子群诱导而来 [1, 3.4.4]，因此上述构造适用。

对每个剩余特征为偶数的局部域，都存在 Weil 群的不可约二维表示不能由真子群诱导而来 [11, 第 29 段]。对于这样的表示，还不知道 $\pi(\sigma)$ 的显式构造。为证明这类 $\pi(\sigma)$ 的存在性，Jacquet 与 Langlands 将局部问题嵌入下面的整体问题。设 $F$ 是整体域（global field），$\rho$ 是 Weil 群 $W_F$ 的不可约连续二维复表示。对 $F$ 的每个素位（place）$v$，令 $\rho_v$ 为 $\rho$ 在局部域 $F_v$ 的 Weil 群上的限制。

\begin{originaltheorem}{定理 1（Theorem 1）[4, 12.2]}
如果对 $\A_F^*/F^*$ 的所有拟特征 $\chi$，整体 $L$ 函数（global $L$-function）$L(s,\rho\otimes\chi)$ 都是复变量 $s$ 的全纯函数，并且在竖直带中有界，则 $\pi(\rho_v)$ 存在。
\end{originaltheorem}

本次会议的基变换（base change）讨论班在讨论 Artin 猜想时，已经介绍了满足定理 1 假设的若干情形。若某个群到 $\GL(2,\C)$ 的同态与到 $\operatorname{PGL}(2,\C)$ 的商映射复合后，其像同构于 $H$，便称这一同态为 $H$ 型（type $H$）。简要地说，定理 1 在下列情形中适用：$\rho$ 由真子群诱导而来（Artin）；$F$ 是正特征域（Weil）；$\rho$ 为 $A_4$ 型（Langlands--Jacquet--Gelbart）；以及 $F=\Q$、$\rho$ 为 $S_4$ 型且复共轭的像的行列式为 $-1$（Langlands--Serre--Deligne）。

\begin{originaltheorem}{定理 2（Theorem 2）[10, 定理 A]}
设 $K$ 是非 Archimedes 局部域，并假设当它为 $\Q_2$ 的真扩域时含有三次单位根。那么对 $\Phi(\GL(2)/K)$ 中所有 $\sigma$，$\pi(\sigma)$ 都存在。
\end{originaltheorem}

证明的关键是构造整体域 $F$、满足 $F_v\simeq K$ 的位 $v$，以及 $W_F$ 的表示 $\rho$，使 $\rho$ 满足定理 1 的假设且 $\rho_v\simeq\sigma$。剩余特征为偶数而不含三次单位根的非 Archimedes 局部域，恰好就是存在 $S_4$ 型 Weil 群表示的那些域 [11, 第 25 段]。对于定理 2 尚未涵盖的域，有些二维 Weil 群表示目前还无法用已知满足定理 1 假设的整体表示来实现。Deligne 指出，如果能够为与全实域（totally real field）$F$ 上除代数（division algebra）所产生的 Shimura 簇（Shimura variety）相关的某些自守表示（automorphic representation）配上 $l$ 进表示（$l$-adic representation），则对于在无穷素位具有指定行为的 $S_4$ 型 $W_F$ 的表示，定理 1 的假设便会成立。于是，当 $K$ 是 $F$ 在整除二的素理想处的完备化（completion）时，定理 2 也就成立；由于 $\Q_2$ 的每个有限扩域都是某个全实域的完备化，第 4 节的命题与定理 3 将由此证明所有情形下的猜想。

\origsection{2}{奇剩余特征（Odd residue characteristic）}
$\GL(2,K)$ 的 Plancherel 公式（Plancherel formula）表明：若 $K$ 的剩余特征为奇数，则形如 $\pi(\sigma)$（$\sigma\in\Phi_{\mathrm{irr}}(\GL(2)/K)$）的超尖点表示穷尽了全部超尖点表示。参考文献 [3]、[8] 和 [9] 论述了奇剩余特征下 $\GL(2,K)$ 及相关群的表示论。

单射性断言可通过考察奇剩余特征情形中超尖点表示的特征标（character）的显式公式来证明；这些公式见 [7] 及上述参考文献（至少对于 $\SL(2)$ 和 $\operatorname{PGL}(2)$ 是如此）。假设 $\sigma$ 由 $K$ 的二次扩域 $E$ 的 Weil 群 $W_E$ 的一维表示 $\lambda$ 诱导而来。对应于 $\lambda$ 的 $E^*$ 的拟特征也用同一符号表示。$\pi(\sigma)$ 的特征标函数限制到一个同构于 $E^*$ 的 Cartan 子群（Cartan subgroup）上，可由一个涉及 $\lambda$ 及其 $\Gal(E/K)$ 共轭的公式给出。由特征标的显式形式可以看出，$\pi(\sigma)$ 在相差 $\Gal(E/K)$ 共轭的意义下确定 $\lambda$，因而确定诱导表示 $\sigma$。

\origsection{3}{正特征（Positive characteristic）}
第 1 节的讨论表明，对于正特征局部域的 Weil 群的所有表示，$\pi(\sigma)$ 都存在。猜想中的双射性断言似乎首先由 Deligne [2] 证明，这是 Drinfeld 关于整体函数域（global function field）上的 $\GL(2)$ 的自守表示与 $l$ 进表示之间关系的结果的一个推论。

\origsection{4}{另一种方法}
[10] 提出了一种方法，为第 2、3 节讨论的情形给出 Langlands 猜想的其他证明，并对定理 2 中的所有域证明了所猜想的双射。

\begin{originaltheorem}{命题（Proposition）[10, 2.2]}
设 $\sigma_1$ 和 $\sigma_2$ 是局部域的 Weil 群的二维表示，且各自都是满足定理 1 假设的某个整体表示的限制。若 $\pi(\sigma_1)\simeq\pi(\sigma_2)$，则 $\sigma_1\simeq\sigma_2$。
\end{originaltheorem}

以上命题通过归纳应用 $\GL(2)$ 的基变换来证明。假设中的条件是必要的，因为局部域上的基变换结果的证明使用了整体方法。

以下结果对于任意非 Archimedes 局部域 $K$ 都成立。

\begin{originaltheorem}{定理 3（Theorem 3）[10, 第 4、5 节]}
可以分别把 $\Phi_{\mathrm{irr}}(\GL(2)/K)$ 和 $\Pi_{\mathrm{cusp}}(\GL(2,K))$ 划分（partition）为有限的子集 $\Phi_\lambda$ 和 $\Pi_\lambda$。两组子集以同一个集合 $\Lambda$ 为指标，并满足：

(1) 若 $\sigma\in\Phi_\lambda$ 且 $\pi(\sigma)$ 存在，则 $\pi(\sigma)\in\Pi_\lambda$。

(2) 对所有 $\lambda\in\Lambda$，均有 $\operatorname{Card}(\Phi_\lambda)=\operatorname{Card}(\Pi_\lambda)$。
\end{originaltheorem}

各划分块由 $\Phi_{\mathrm{irr}}(\GL(2)/K)$ 中元素的 Artin 导子（Artin conductor）和行列式上的条件（分别对应 $\Pi_{\mathrm{cusp}}(\GL(2,K))$ 中元素的导子与中心拟特征上的条件）确定。由于表示的导子由扭曲 $\varepsilon$ 因子确定，第一项断言直接来自定义。计算 $\Phi_\lambda$ 的基数时，先依照 [11] 构造 $W_K$ 的所有不可约二维表示，再统计其中具有指定 Artin 导子和行列式的表示。集合 $\Pi_\lambda$ 则通过 $\GL(2,K)$ 的平方可积表示（square-integrable representation），与 $K$ 上四元数除代数的可逆元群的容许表示之间的对应来研究。

定理 2 加上单射性命题和定理 3 的计数结果，表明对于满足定理 2 假设的局部域，$\Phi_\lambda$ 与 $\Pi_\lambda$ 通过 $\sigma\mapsto\pi(\sigma)$ 双射对应。由此便得到引言中列出的各个已知情形。

\origsection{5}{注记（Remarks）}
虽然 Langlands 猜想的陈述纯粹是局部的，但上述偶剩余特征情形的证明使用了整体方法（基变换和 Artin 猜想的若干情形）。只有奇剩余特征情形的上述证明才是纯局部的。

Cartier 与 Nobs 提出了域 $\Q_2$ 上这一猜想的一个纯局部证明。他们计算 $W_{\Q_2}$ 的不可约二维表示所需的 $\varepsilon$ 因子，并把它们与 [6] 构造的 $\GL(2,\Q_2)$ 的超尖点表示的因子匹配。这些超尖点表示通过诱导 $\GL(2,K)$ 的模中心紧（compact modulo the center）的子群的有限维表示来构造。[5] 给出了 $\GL(2,K)$ 的超尖点表示的一种类似构造；原则上，借助这一构造可以在一般情形下计算 $\varepsilon$ 因子，并与 $W_K$ 表示的相应因子进行匹配。

\begin{thebibliography}{11}
\bibitem[1]{c28-1} P. Deligne, \emph{Formes modulaires et représentations de GL(2)}, in \emph{Modular functions of one variable}, Springer Lecture Notes, No. 349, 1973, pp. 55--106.
\bibitem[2]{c28-2} P. Deligne, \emph{Review of ``Elliptic A-modules'' by V. Drinfeld}, Math. Reviews 52 (1976), \#5580.
\bibitem[3]{c28-3} I. M. Gelfand, M. I. Graev and I. I. Piatetskii-Shapiro, \emph{Representation theory and automorphic functions}, Saunders, 1969.
\bibitem[4]{c28-4} H. Jacquet and R. P. Langlands, \emph{Automorphic forms on GL(2)}, Springer Lecture Notes, No. 114, 1970.
\bibitem[5]{c28-5} P. Kutzko, \emph{On the supercuspidal representations of $\GL_2$. I, II}, Amer. J. Math. 100 (1978), 43--60.
\bibitem[6]{c28-6} A. Nobs, \emph{Supercuspidal representations of $\GL(2,\Q_p)$, including $p=2$}（待刊）。
\bibitem[7]{c28-7} P. J. Sally, Jr. and J. A. Shalika, \emph{Characters of the discrete series of representations of SL(2) over a local field}, Proc. Nat. Acad. Sci. U.S.A. 61 (1968), 1231--1237.
\bibitem[8]{c28-8} P. J. Sally, Jr. and J. A. Shalika, \emph{The Plancherel formula for SL(2) over a local field}, Proc. Nat. Acad. Sci. U.S.A. 63 (1969), 661--667.
\bibitem[9]{c28-9} A. J. Silberger, \emph{PGL(2) over the $p$-adics}, Springer Lecture Notes, No. 166, 1970.
\bibitem[10]{c28-10} J. Tunnell, \emph{On the local Langlands Conjecture for GL(2)}, Thesis, Harvard University (1977).
\bibitem[11]{c28-11} A. Weil, \emph{Exercices dyadiques}, Invent. Math. 27 (1974), 1--22.
\end{thebibliography}
\begin{flushleft}\small Princeton University\end{flushleft}
